\scanpage{10}
satisfaisant à la propriété d'associativité évidente, par prolongement par continuité à partir de l'accouplement au niveau des $\operatorname{Hom}$ sans $\hat{\ }$. \deleted{La construction marche dans le \unsure{cas}}

\margin{Mérite d'être faite avec plus de soin !}
\uline{Digression axiomatique}\quad Une catégorie $\unsure{\mathcal{C}|C}$ est dite \deleted{``à $\operatorname{Hom}$ profinis'' si les $\operatorname{Hom}$} \deleted{``à $\operatorname{Hom}$''} ``à \deleted{autom} \inserted{groupes d'autom.} profinis'' si pour tt $X\in C$, $\operatorname{Aut}X$ est muni d'une structure de groupe profini, de telle façon que pour tout $u:X\longrightarrow Y$ de $C$, les stabilisateurs de $u\in\operatorname{Hom}_C(X,Y)$ dans $\operatorname{Aut}_C(X)\times\operatorname{Aut}_C(Y)$ \deleted{\unsure{à}} \inserted{opérant par $u\longmapsto guf^{-1}$,} \deleted{droite) et $\operatorname{Aut}_C(Y)$ (opère à g.) sont des} sous-groupes fermés. \deleted{\unsure{On}} On munit alors chaque orbite de $\operatorname{Hom}(X,Y)$ sous $\operatorname{Aut}(X)\times\operatorname{Aut}(Y)$ de la topologie profinie quotient, et $\operatorname{Hom}(X,Y)$ de la topologie \inserted{disjointe} somme des orbites, de sorte que $\operatorname{Hom}(X,Y)$ devient un espace loc. compact.

Je \deleted{\unsure{dis}} \inserted{veux} que l'accouplement des morphismes
\[
(14)\quad\left\{
\begin{aligned}
\operatorname{Hom}(X,Y)\times\operatorname{Hom}(Y,Z)&\mathbin{\unsure{-|\longrightarrow}}\operatorname{Hom}(X,Z)\\
(u,v)&\longmapsto vu
\end{aligned}\right.
\]
est alors \deleted{continue} une \deleted{\unsure{morphi.}} application continue. Il suffit de le voir sur des produits d'orbites, donc de voir que \inserted{pour $u,v$ fixés,} l'application
\[
\begin{gathered}
(g,f,g',f')\longmapsto(g'vf'^{-1})\circ(guf^{-1})\\
\operatorname{Aut}(\unsure{Y|Z})\times\operatorname{Aut}(X)\times\operatorname{Aut}(Z)\times\operatorname{Aut}(Y)\longrightarrow\operatorname{Hom}(X,Z)
\end{gathered}
\]
\margin{9}
